% explainer-source-hash: 203ccc1f27de4a9304ff1611711b63dc30f9e2bab3fc84a81cc3da8ca4b298c1
% explainer-source-commit: a254e39de5ae4916966daf41ae0ea1353e3e4ea3
% explainer-generated: 2026-10-27
% Regenerate with /explain-all (see WIRING.md). Do not edit this stamp by hand:
% it is how the toolchain knows whether this document still matches the code.
\documentclass[11pt,a4paper]{article}
% preamble.tex - shared preamble for this project's explainer documents.
%
% This file is generated: it is copied in beside the documents on every build, so a
% clone of this repo can rebuild the PDFs, but local edits to it are overwritten.
%
% Usage from a document:
%   \documentclass[11pt,a4paper]{article}
%   % preamble.tex - shared preamble for this project's explainer documents.
%
% This file is generated: it is copied in beside the documents on every build, so a
% clone of this repo can rebuild the PDFs, but local edits to it are overwritten.
%
% Usage from a document:
%   \documentclass[11pt,a4paper]{article}
%   % preamble.tex - shared preamble for this project's explainer documents.
%
% This file is generated: it is copied in beside the documents on every build, so a
% clone of this repo can rebuild the PDFs, but local edits to it are overwritten.
%
% Usage from a document:
%   \documentclass[11pt,a4paper]{article}
%   \input{preamble}
%   \begin{document} ... \end{document}
%
% Build:  pdflatex -interaction=nonstopmode deep-dive.tex   (run it three times, so the
%         table of contents and the TikZ positioning settle)
%
% PACKAGE NOTES. Every package used here was checked present on the machine these
% documents were built on (Debian TeX Live 2023). If you rebuild elsewhere and a build
% fails, these are the likely gaps. Do not add a package without checking first:
%   kpsewhich <pkg>.sty
%
% Deliberately NOT used, because they were absent and the build dies without them:
%   algorithm2e, algorithm, algorithmic, algpseudocode  (texlive-science not installed)
%     -> the \begin{pseudocode} environment defined below stands in for them.
%   lmodern, charter, mathptmx, times, mathpazo, courier, newtx*, tgtermes
%     -> their .sty files RESOLVE but the font files themselves are absent, so the
%        build dies with "I can't find file `bchr8t'" or similar. Verified by
%        test-compiling each one. Only Computer Modern (the default) and helvet
%        actually render. kpsewhich reporting a font it cannot render is the trap:
%        do not "upgrade" the font stack without test-compiling it first.
%   amsmath, amssymb -> not loaded, so \text{} and \blacksquare are unavailable.
%
% A fuller font stack and algorithm2e would come from:
%   sudo apt install texlive-fonts-recommended texlive-science cm-super

\usepackage[T1]{fontenc}
\usepackage[utf8]{inputenc}
% Font: Computer Modern (default). See the note above before changing this.
\usepackage[a4paper,margin=2.4cm,headheight=14pt]{geometry}
% microtype WITHOUT font expansion. This is load-bearing, do not "restore" it:
% cm-super is not installed, so T1 Computer Modern resolves to BITMAP fonts, and
% pdfTeX's auto-expansion needs scalable ones. With expansion on, any document
% long enough to reach a second page dies with
%   "pdfTeX error (font expansion): auto expansion is only possible with scalable fonts"
% Protrusion (the bigger visual win) still works and stays on.
% `sudo apt install cm-super` would lift this; until then, leave expansion off.
\usepackage[protrusion=true,expansion=false]{microtype}
% Kill hyphen ligatures in TYPEWRITER text. This is load-bearing, not cosmetic. Under T1,
% \texttt{--verbose} silently renders as "-verbose", and \texttt{---} as "--", so a
% document can print a command-line flag as the WRONG flag and look perfectly fine doing
% it. Verified by rendering the page and reading the glyphs. lstlisting and \verb were
% always safe; this makes \code, \file and \texttt safe too, which is what prose uses.
% Ordinary prose is unaffected and still ligates, which is what you want there.
\DisableLigatures[-]{family=tt*}
\usepackage{graphicx}
\usepackage{xcolor}
\usepackage{booktabs}
\usepackage{enumitem}
\usepackage{fancyhdr}
\usepackage{titlesec}
\usepackage{tcolorbox}
\usepackage{listings}
\usepackage{float}
\usepackage{caption}
\usepackage{tikz}
\usepackage{forest}
\usepackage{pgfplots}
\pgfplotsset{compat=1.18}
\usetikzlibrary{arrows.meta,positioning,shapes.geometric,fit,backgrounds,calc,decorations.pathreplacing}
\tcbuselibrary{skins,breakable}
\usepackage[hidelinks,bookmarks=true,bookmarksnumbered=true]{hyperref}

% ===== palette ===========================================================
\definecolor{inkblue}{HTML}{1F3A5F}
\definecolor{accent}{HTML}{2E7DAF}
\definecolor{softgrey}{HTML}{F4F6F8}
\definecolor{linegrey}{HTML}{C8D0D8}
\definecolor{codebg}{HTML}{F7F7F4}
\definecolor{warnred}{HTML}{A33B3B}
\definecolor{okgreen}{HTML}{2E7D53}

% ===== headings ==========================================================
\titleformat{\section}{\Large\bfseries\color{inkblue}}{\thesection}{0.7em}{}
\titleformat{\subsection}{\large\bfseries\color{inkblue}}{\thesubsection}{0.6em}{}
\titleformat{\subsubsection}{\normalsize\bfseries\color{accent}}{\thesubsubsection}{0.5em}{}

\pagestyle{fancy}
\fancyhf{}
\fancyhead[L]{\small\color{inkblue}\nouppercase{\leftmark}}
\fancyhead[R]{\small\color{inkblue}\thepage}
\renewcommand{\headrulewidth}{0.4pt}
\renewcommand{\headrule}{\hbox to\headwidth{\color{linegrey}\leaders\hrule height \headrulewidth\hfill}}

% ===== code ==============================================================
\lstdefinestyle{house}{
  backgroundcolor=\color{codebg},
  basicstyle=\ttfamily\footnotesize,
  keywordstyle=\color{inkblue}\bfseries,
  commentstyle=\color{gray}\itshape,
  stringstyle=\color{okgreen},
  numbers=left, numberstyle=\tiny\color{gray}, numbersep=8pt,
  frame=single, rulecolor=\color{linegrey},
  breaklines=true, breakatwhitespace=false,
  showstringspaces=false, tabsize=2,
  captionpos=b, xleftmargin=14pt, framexleftmargin=14pt,
  columns=fullflexible, keepspaces=true,
  upquote=true,  % render ' and " as straight quotes, not curly
  % NON-ASCII INSIDE A LISTING IS A HARD BUILD FAILURE, not a cosmetic issue: listings
  % reads raw bytes, so inputenc dies on the first multi-byte character ("Invalid UTF-8
  % byte sequence") and the whole document fails to build. These mappings cover what
  % realistically turns up in source code: Latin-1 accents, Spanish punctuation, smart
  % quotes picked up from a word processor, the degree and section signs.
  % NOT COVERED, and it cannot be here: Urdu, Arabic and Balochi script. pdflatex has no
  % font for them. Describe or transliterate such text; never paste it into a listing.
  literate=%
    {á}{{\'a}}1 {é}{{\'e}}1 {í}{{\'i}}1 {ó}{{\'o}}1 {ú}{{\'u}}1
    {Á}{{\'A}}1 {É}{{\'E}}1 {Í}{{\'I}}1 {Ó}{{\'O}}1 {Ú}{{\'U}}1
    {à}{{\`a}}1 {è}{{\`e}}1 {ì}{{\`i}}1 {ò}{{\`o}}1 {ù}{{\`u}}1
    {ä}{{\"a}}1 {ë}{{\"e}}1 {ï}{{\"i}}1 {ö}{{\"o}}1 {ü}{{\"u}}1
    {Ä}{{\"A}}1 {Ö}{{\"O}}1 {Ü}{{\"U}}1 {ß}{{\ss}}1
    {ñ}{{\~n}}1 {Ñ}{{\~N}}1 {ç}{{\c c}}1 {Ç}{{\c C}}1
    {¿}{{\textquestiondown}}1 {¡}{{\textexclamdown}}1
    {°}{{\textdegree}}1 {§}{{\S}}1 {£}{{\pounds}}1 {€}{{EUR}}1
    {“}{{\textquotedblleft}}1 {”}{{\textquotedblright}}1
    {‘}{{\textquoteleft}}1 {’}{{\textquoteright}}1
    {–}{{-}}1 {—}{{-}}1 {…}{{...}}1 {→}{{$\rightarrow$}}1 {✓}{{x}}1
}
\lstset{style=house}

% ===== callout boxes =====================================================
% \begin{keyidea}{Title} ... \end{keyidea}   - a concept a beginner must grasp
\newtcolorbox{keyidea}[1]{
  enhanced, breakable, colback=softgrey, colframe=accent,
  boxrule=0.4pt, arc=2pt, left=8pt, right=8pt, top=6pt, bottom=6pt,
  title={\bfseries #1}, coltitle=white, colbacktitle=accent, fonttitle=\small
}

% \begin{jargon}{Term} definition \end{jargon}  - define a term at first use
\newtcolorbox{jargon}[1]{
  enhanced, breakable, colback=white, colframe=linegrey,
  boxrule=0.4pt, arc=2pt, left=8pt, right=8pt, top=6pt, bottom=6pt,
  title={\bfseries Jargon: #1}, coltitle=inkblue, colbacktitle=softgrey, fonttitle=\small
}

% \begin{honest}{Title} ... \end{honest}  - brutally honest finding (rule 18)
\newtcolorbox{honest}[1]{
  enhanced, breakable, colback=white, colframe=warnred,
  boxrule=0.6pt, arc=2pt, left=8pt, right=8pt, top=6pt, bottom=6pt,
  title={\bfseries #1}, coltitle=white, colbacktitle=warnred, fonttitle=\small
}

% \begin{pseudocode}{Caption} ... \end{pseudocode}
% Stands in for algorithm2e, which is NOT installed. Use \Step / \Indent inside.
\newtcolorbox{pseudocode}[1]{
  enhanced, breakable, colback=codebg, colframe=inkblue,
  boxrule=0.4pt, arc=2pt, left=10pt, right=8pt, top=6pt, bottom=6pt,
  fontupper=\ttfamily\small,
  title={\bfseries\rmfamily #1}, coltitle=white, colbacktitle=inkblue, fonttitle=\small
}
\newcommand{\Step}[1]{\par\noindent #1}
\newcommand{\Indent}[1]{\par\noindent\hspace*{2em}#1}

% ===== tikz house styles =================================================
\tikzset{
  box/.style   = {rectangle, rounded corners=2pt, draw=inkblue, fill=softgrey,
                  text=inkblue, align=center, inner sep=6pt, minimum height=9mm,
                  font=\small},
  extbox/.style= {rectangle, rounded corners=2pt, draw=linegrey, fill=white,
                  text=black!70, align=center, inner sep=6pt, minimum height=9mm,
                  font=\small, dashed},
  store/.style = {cylinder, shape border rotate=90, aspect=0.25, draw=accent,
                  fill=white, text=inkblue, align=center, inner sep=5pt,
                  font=\small},
  dec/.style   = {diamond, aspect=1.6, draw=accent, fill=white, text=inkblue,
                  align=center, inner sep=2pt, font=\scriptsize},
  flow/.style  = {-{Stealth[length=2.4mm]}, draw=inkblue, thick},
  dflow/.style = {-{Stealth[length=2.4mm]}, draw=accent, thick, dashed},
  % align=center is load-bearing: without it, a \\ inside an edge label throws
  % the misleading "perhaps a missing \item" error.
  lbl/.style   = {font=\scriptsize\itshape, text=black!65, fill=white,
                  inner sep=2pt, align=center}
}

% forest default for directory trees
\forestset{
  dirtree/.style={
    for tree={
      font=\ttfamily\small, grow'=0, child anchor=west, parent anchor=south,
      anchor=west, calign=first, inner sep=2pt,
      edge path={\noexpand\path [draw, \forestoption{edge}]
        (!u.south west) +(9pt,0) |- (.child anchor) \forestoption{edge label};},
      before typesetting nodes={if n=1{insert before={[,phantom]}}{}},
      fit=band, before computing xy={l=15pt}
    }
  }
}

% ===== misc ==============================================================
\captionsetup{font=small, labelfont={bf,color=inkblue}}
\setlist[itemize]{leftmargin=1.2em, itemsep=2pt, topsep=3pt}
\setlist[enumerate]{leftmargin=1.4em, itemsep=2pt, topsep=3pt}
\newcommand{\file}[1]{\texttt{\small #1}}
\newcommand{\code}[1]{\texttt{\small #1}}
\setcounter{tocdepth}{2}
\setcounter{secnumdepth}{3}

%   \begin{document} ... \end{document}
%
% Build:  pdflatex -interaction=nonstopmode deep-dive.tex   (run it three times, so the
%         table of contents and the TikZ positioning settle)
%
% PACKAGE NOTES. Every package used here was checked present on the machine these
% documents were built on (Debian TeX Live 2023). If you rebuild elsewhere and a build
% fails, these are the likely gaps. Do not add a package without checking first:
%   kpsewhich <pkg>.sty
%
% Deliberately NOT used, because they were absent and the build dies without them:
%   algorithm2e, algorithm, algorithmic, algpseudocode  (texlive-science not installed)
%     -> the \begin{pseudocode} environment defined below stands in for them.
%   lmodern, charter, mathptmx, times, mathpazo, courier, newtx*, tgtermes
%     -> their .sty files RESOLVE but the font files themselves are absent, so the
%        build dies with "I can't find file `bchr8t'" or similar. Verified by
%        test-compiling each one. Only Computer Modern (the default) and helvet
%        actually render. kpsewhich reporting a font it cannot render is the trap:
%        do not "upgrade" the font stack without test-compiling it first.
%   amsmath, amssymb -> not loaded, so \text{} and \blacksquare are unavailable.
%
% A fuller font stack and algorithm2e would come from:
%   sudo apt install texlive-fonts-recommended texlive-science cm-super

\usepackage[T1]{fontenc}
\usepackage[utf8]{inputenc}
% Font: Computer Modern (default). See the note above before changing this.
\usepackage[a4paper,margin=2.4cm,headheight=14pt]{geometry}
% microtype WITHOUT font expansion. This is load-bearing, do not "restore" it:
% cm-super is not installed, so T1 Computer Modern resolves to BITMAP fonts, and
% pdfTeX's auto-expansion needs scalable ones. With expansion on, any document
% long enough to reach a second page dies with
%   "pdfTeX error (font expansion): auto expansion is only possible with scalable fonts"
% Protrusion (the bigger visual win) still works and stays on.
% `sudo apt install cm-super` would lift this; until then, leave expansion off.
\usepackage[protrusion=true,expansion=false]{microtype}
% Kill hyphen ligatures in TYPEWRITER text. This is load-bearing, not cosmetic. Under T1,
% \texttt{--verbose} silently renders as "-verbose", and \texttt{---} as "--", so a
% document can print a command-line flag as the WRONG flag and look perfectly fine doing
% it. Verified by rendering the page and reading the glyphs. lstlisting and \verb were
% always safe; this makes \code, \file and \texttt safe too, which is what prose uses.
% Ordinary prose is unaffected and still ligates, which is what you want there.
\DisableLigatures[-]{family=tt*}
\usepackage{graphicx}
\usepackage{xcolor}
\usepackage{booktabs}
\usepackage{enumitem}
\usepackage{fancyhdr}
\usepackage{titlesec}
\usepackage{tcolorbox}
\usepackage{listings}
\usepackage{float}
\usepackage{caption}
\usepackage{tikz}
\usepackage{forest}
\usepackage{pgfplots}
\pgfplotsset{compat=1.18}
\usetikzlibrary{arrows.meta,positioning,shapes.geometric,fit,backgrounds,calc,decorations.pathreplacing}
\tcbuselibrary{skins,breakable}
\usepackage[hidelinks,bookmarks=true,bookmarksnumbered=true]{hyperref}

% ===== palette ===========================================================
\definecolor{inkblue}{HTML}{1F3A5F}
\definecolor{accent}{HTML}{2E7DAF}
\definecolor{softgrey}{HTML}{F4F6F8}
\definecolor{linegrey}{HTML}{C8D0D8}
\definecolor{codebg}{HTML}{F7F7F4}
\definecolor{warnred}{HTML}{A33B3B}
\definecolor{okgreen}{HTML}{2E7D53}

% ===== headings ==========================================================
\titleformat{\section}{\Large\bfseries\color{inkblue}}{\thesection}{0.7em}{}
\titleformat{\subsection}{\large\bfseries\color{inkblue}}{\thesubsection}{0.6em}{}
\titleformat{\subsubsection}{\normalsize\bfseries\color{accent}}{\thesubsubsection}{0.5em}{}

\pagestyle{fancy}
\fancyhf{}
\fancyhead[L]{\small\color{inkblue}\nouppercase{\leftmark}}
\fancyhead[R]{\small\color{inkblue}\thepage}
\renewcommand{\headrulewidth}{0.4pt}
\renewcommand{\headrule}{\hbox to\headwidth{\color{linegrey}\leaders\hrule height \headrulewidth\hfill}}

% ===== code ==============================================================
\lstdefinestyle{house}{
  backgroundcolor=\color{codebg},
  basicstyle=\ttfamily\footnotesize,
  keywordstyle=\color{inkblue}\bfseries,
  commentstyle=\color{gray}\itshape,
  stringstyle=\color{okgreen},
  numbers=left, numberstyle=\tiny\color{gray}, numbersep=8pt,
  frame=single, rulecolor=\color{linegrey},
  breaklines=true, breakatwhitespace=false,
  showstringspaces=false, tabsize=2,
  captionpos=b, xleftmargin=14pt, framexleftmargin=14pt,
  columns=fullflexible, keepspaces=true,
  upquote=true,  % render ' and " as straight quotes, not curly
  % NON-ASCII INSIDE A LISTING IS A HARD BUILD FAILURE, not a cosmetic issue: listings
  % reads raw bytes, so inputenc dies on the first multi-byte character ("Invalid UTF-8
  % byte sequence") and the whole document fails to build. These mappings cover what
  % realistically turns up in source code: Latin-1 accents, Spanish punctuation, smart
  % quotes picked up from a word processor, the degree and section signs.
  % NOT COVERED, and it cannot be here: Urdu, Arabic and Balochi script. pdflatex has no
  % font for them. Describe or transliterate such text; never paste it into a listing.
  literate=%
    {á}{{\'a}}1 {é}{{\'e}}1 {í}{{\'i}}1 {ó}{{\'o}}1 {ú}{{\'u}}1
    {Á}{{\'A}}1 {É}{{\'E}}1 {Í}{{\'I}}1 {Ó}{{\'O}}1 {Ú}{{\'U}}1
    {à}{{\`a}}1 {è}{{\`e}}1 {ì}{{\`i}}1 {ò}{{\`o}}1 {ù}{{\`u}}1
    {ä}{{\"a}}1 {ë}{{\"e}}1 {ï}{{\"i}}1 {ö}{{\"o}}1 {ü}{{\"u}}1
    {Ä}{{\"A}}1 {Ö}{{\"O}}1 {Ü}{{\"U}}1 {ß}{{\ss}}1
    {ñ}{{\~n}}1 {Ñ}{{\~N}}1 {ç}{{\c c}}1 {Ç}{{\c C}}1
    {¿}{{\textquestiondown}}1 {¡}{{\textexclamdown}}1
    {°}{{\textdegree}}1 {§}{{\S}}1 {£}{{\pounds}}1 {€}{{EUR}}1
    {“}{{\textquotedblleft}}1 {”}{{\textquotedblright}}1
    {‘}{{\textquoteleft}}1 {’}{{\textquoteright}}1
    {–}{{-}}1 {—}{{-}}1 {…}{{...}}1 {→}{{$\rightarrow$}}1 {✓}{{x}}1
}
\lstset{style=house}

% ===== callout boxes =====================================================
% \begin{keyidea}{Title} ... \end{keyidea}   - a concept a beginner must grasp
\newtcolorbox{keyidea}[1]{
  enhanced, breakable, colback=softgrey, colframe=accent,
  boxrule=0.4pt, arc=2pt, left=8pt, right=8pt, top=6pt, bottom=6pt,
  title={\bfseries #1}, coltitle=white, colbacktitle=accent, fonttitle=\small
}

% \begin{jargon}{Term} definition \end{jargon}  - define a term at first use
\newtcolorbox{jargon}[1]{
  enhanced, breakable, colback=white, colframe=linegrey,
  boxrule=0.4pt, arc=2pt, left=8pt, right=8pt, top=6pt, bottom=6pt,
  title={\bfseries Jargon: #1}, coltitle=inkblue, colbacktitle=softgrey, fonttitle=\small
}

% \begin{honest}{Title} ... \end{honest}  - brutally honest finding (rule 18)
\newtcolorbox{honest}[1]{
  enhanced, breakable, colback=white, colframe=warnred,
  boxrule=0.6pt, arc=2pt, left=8pt, right=8pt, top=6pt, bottom=6pt,
  title={\bfseries #1}, coltitle=white, colbacktitle=warnred, fonttitle=\small
}

% \begin{pseudocode}{Caption} ... \end{pseudocode}
% Stands in for algorithm2e, which is NOT installed. Use \Step / \Indent inside.
\newtcolorbox{pseudocode}[1]{
  enhanced, breakable, colback=codebg, colframe=inkblue,
  boxrule=0.4pt, arc=2pt, left=10pt, right=8pt, top=6pt, bottom=6pt,
  fontupper=\ttfamily\small,
  title={\bfseries\rmfamily #1}, coltitle=white, colbacktitle=inkblue, fonttitle=\small
}
\newcommand{\Step}[1]{\par\noindent #1}
\newcommand{\Indent}[1]{\par\noindent\hspace*{2em}#1}

% ===== tikz house styles =================================================
\tikzset{
  box/.style   = {rectangle, rounded corners=2pt, draw=inkblue, fill=softgrey,
                  text=inkblue, align=center, inner sep=6pt, minimum height=9mm,
                  font=\small},
  extbox/.style= {rectangle, rounded corners=2pt, draw=linegrey, fill=white,
                  text=black!70, align=center, inner sep=6pt, minimum height=9mm,
                  font=\small, dashed},
  store/.style = {cylinder, shape border rotate=90, aspect=0.25, draw=accent,
                  fill=white, text=inkblue, align=center, inner sep=5pt,
                  font=\small},
  dec/.style   = {diamond, aspect=1.6, draw=accent, fill=white, text=inkblue,
                  align=center, inner sep=2pt, font=\scriptsize},
  flow/.style  = {-{Stealth[length=2.4mm]}, draw=inkblue, thick},
  dflow/.style = {-{Stealth[length=2.4mm]}, draw=accent, thick, dashed},
  % align=center is load-bearing: without it, a \\ inside an edge label throws
  % the misleading "perhaps a missing \item" error.
  lbl/.style   = {font=\scriptsize\itshape, text=black!65, fill=white,
                  inner sep=2pt, align=center}
}

% forest default for directory trees
\forestset{
  dirtree/.style={
    for tree={
      font=\ttfamily\small, grow'=0, child anchor=west, parent anchor=south,
      anchor=west, calign=first, inner sep=2pt,
      edge path={\noexpand\path [draw, \forestoption{edge}]
        (!u.south west) +(9pt,0) |- (.child anchor) \forestoption{edge label};},
      before typesetting nodes={if n=1{insert before={[,phantom]}}{}},
      fit=band, before computing xy={l=15pt}
    }
  }
}

% ===== misc ==============================================================
\captionsetup{font=small, labelfont={bf,color=inkblue}}
\setlist[itemize]{leftmargin=1.2em, itemsep=2pt, topsep=3pt}
\setlist[enumerate]{leftmargin=1.4em, itemsep=2pt, topsep=3pt}
\newcommand{\file}[1]{\texttt{\small #1}}
\newcommand{\code}[1]{\texttt{\small #1}}
\setcounter{tocdepth}{2}
\setcounter{secnumdepth}{3}

%   \begin{document} ... \end{document}
%
% Build:  pdflatex -interaction=nonstopmode deep-dive.tex   (run it three times, so the
%         table of contents and the TikZ positioning settle)
%
% PACKAGE NOTES. Every package used here was checked present on the machine these
% documents were built on (Debian TeX Live 2023). If you rebuild elsewhere and a build
% fails, these are the likely gaps. Do not add a package without checking first:
%   kpsewhich <pkg>.sty
%
% Deliberately NOT used, because they were absent and the build dies without them:
%   algorithm2e, algorithm, algorithmic, algpseudocode  (texlive-science not installed)
%     -> the \begin{pseudocode} environment defined below stands in for them.
%   lmodern, charter, mathptmx, times, mathpazo, courier, newtx*, tgtermes
%     -> their .sty files RESOLVE but the font files themselves are absent, so the
%        build dies with "I can't find file `bchr8t'" or similar. Verified by
%        test-compiling each one. Only Computer Modern (the default) and helvet
%        actually render. kpsewhich reporting a font it cannot render is the trap:
%        do not "upgrade" the font stack without test-compiling it first.
%   amsmath, amssymb -> not loaded, so \text{} and \blacksquare are unavailable.
%
% A fuller font stack and algorithm2e would come from:
%   sudo apt install texlive-fonts-recommended texlive-science cm-super

\usepackage[T1]{fontenc}
\usepackage[utf8]{inputenc}
% Font: Computer Modern (default). See the note above before changing this.
\usepackage[a4paper,margin=2.4cm,headheight=14pt]{geometry}
% microtype WITHOUT font expansion. This is load-bearing, do not "restore" it:
% cm-super is not installed, so T1 Computer Modern resolves to BITMAP fonts, and
% pdfTeX's auto-expansion needs scalable ones. With expansion on, any document
% long enough to reach a second page dies with
%   "pdfTeX error (font expansion): auto expansion is only possible with scalable fonts"
% Protrusion (the bigger visual win) still works and stays on.
% `sudo apt install cm-super` would lift this; until then, leave expansion off.
\usepackage[protrusion=true,expansion=false]{microtype}
% Kill hyphen ligatures in TYPEWRITER text. This is load-bearing, not cosmetic. Under T1,
% \texttt{--verbose} silently renders as "-verbose", and \texttt{---} as "--", so a
% document can print a command-line flag as the WRONG flag and look perfectly fine doing
% it. Verified by rendering the page and reading the glyphs. lstlisting and \verb were
% always safe; this makes \code, \file and \texttt safe too, which is what prose uses.
% Ordinary prose is unaffected and still ligates, which is what you want there.
\DisableLigatures[-]{family=tt*}
\usepackage{graphicx}
\usepackage{xcolor}
\usepackage{booktabs}
\usepackage{enumitem}
\usepackage{fancyhdr}
\usepackage{titlesec}
\usepackage{tcolorbox}
\usepackage{listings}
\usepackage{float}
\usepackage{caption}
\usepackage{tikz}
\usepackage{forest}
\usepackage{pgfplots}
\pgfplotsset{compat=1.18}
\usetikzlibrary{arrows.meta,positioning,shapes.geometric,fit,backgrounds,calc,decorations.pathreplacing}
\tcbuselibrary{skins,breakable}
\usepackage[hidelinks,bookmarks=true,bookmarksnumbered=true]{hyperref}

% ===== palette ===========================================================
\definecolor{inkblue}{HTML}{1F3A5F}
\definecolor{accent}{HTML}{2E7DAF}
\definecolor{softgrey}{HTML}{F4F6F8}
\definecolor{linegrey}{HTML}{C8D0D8}
\definecolor{codebg}{HTML}{F7F7F4}
\definecolor{warnred}{HTML}{A33B3B}
\definecolor{okgreen}{HTML}{2E7D53}

% ===== headings ==========================================================
\titleformat{\section}{\Large\bfseries\color{inkblue}}{\thesection}{0.7em}{}
\titleformat{\subsection}{\large\bfseries\color{inkblue}}{\thesubsection}{0.6em}{}
\titleformat{\subsubsection}{\normalsize\bfseries\color{accent}}{\thesubsubsection}{0.5em}{}

\pagestyle{fancy}
\fancyhf{}
\fancyhead[L]{\small\color{inkblue}\nouppercase{\leftmark}}
\fancyhead[R]{\small\color{inkblue}\thepage}
\renewcommand{\headrulewidth}{0.4pt}
\renewcommand{\headrule}{\hbox to\headwidth{\color{linegrey}\leaders\hrule height \headrulewidth\hfill}}

% ===== code ==============================================================
\lstdefinestyle{house}{
  backgroundcolor=\color{codebg},
  basicstyle=\ttfamily\footnotesize,
  keywordstyle=\color{inkblue}\bfseries,
  commentstyle=\color{gray}\itshape,
  stringstyle=\color{okgreen},
  numbers=left, numberstyle=\tiny\color{gray}, numbersep=8pt,
  frame=single, rulecolor=\color{linegrey},
  breaklines=true, breakatwhitespace=false,
  showstringspaces=false, tabsize=2,
  captionpos=b, xleftmargin=14pt, framexleftmargin=14pt,
  columns=fullflexible, keepspaces=true,
  upquote=true,  % render ' and " as straight quotes, not curly
  % NON-ASCII INSIDE A LISTING IS A HARD BUILD FAILURE, not a cosmetic issue: listings
  % reads raw bytes, so inputenc dies on the first multi-byte character ("Invalid UTF-8
  % byte sequence") and the whole document fails to build. These mappings cover what
  % realistically turns up in source code: Latin-1 accents, Spanish punctuation, smart
  % quotes picked up from a word processor, the degree and section signs.
  % NOT COVERED, and it cannot be here: Urdu, Arabic and Balochi script. pdflatex has no
  % font for them. Describe or transliterate such text; never paste it into a listing.
  literate=%
    {á}{{\'a}}1 {é}{{\'e}}1 {í}{{\'i}}1 {ó}{{\'o}}1 {ú}{{\'u}}1
    {Á}{{\'A}}1 {É}{{\'E}}1 {Í}{{\'I}}1 {Ó}{{\'O}}1 {Ú}{{\'U}}1
    {à}{{\`a}}1 {è}{{\`e}}1 {ì}{{\`i}}1 {ò}{{\`o}}1 {ù}{{\`u}}1
    {ä}{{\"a}}1 {ë}{{\"e}}1 {ï}{{\"i}}1 {ö}{{\"o}}1 {ü}{{\"u}}1
    {Ä}{{\"A}}1 {Ö}{{\"O}}1 {Ü}{{\"U}}1 {ß}{{\ss}}1
    {ñ}{{\~n}}1 {Ñ}{{\~N}}1 {ç}{{\c c}}1 {Ç}{{\c C}}1
    {¿}{{\textquestiondown}}1 {¡}{{\textexclamdown}}1
    {°}{{\textdegree}}1 {§}{{\S}}1 {£}{{\pounds}}1 {€}{{EUR}}1
    {“}{{\textquotedblleft}}1 {”}{{\textquotedblright}}1
    {‘}{{\textquoteleft}}1 {’}{{\textquoteright}}1
    {–}{{-}}1 {—}{{-}}1 {…}{{...}}1 {→}{{$\rightarrow$}}1 {✓}{{x}}1
}
\lstset{style=house}

% ===== callout boxes =====================================================
% \begin{keyidea}{Title} ... \end{keyidea}   - a concept a beginner must grasp
\newtcolorbox{keyidea}[1]{
  enhanced, breakable, colback=softgrey, colframe=accent,
  boxrule=0.4pt, arc=2pt, left=8pt, right=8pt, top=6pt, bottom=6pt,
  title={\bfseries #1}, coltitle=white, colbacktitle=accent, fonttitle=\small
}

% \begin{jargon}{Term} definition \end{jargon}  - define a term at first use
\newtcolorbox{jargon}[1]{
  enhanced, breakable, colback=white, colframe=linegrey,
  boxrule=0.4pt, arc=2pt, left=8pt, right=8pt, top=6pt, bottom=6pt,
  title={\bfseries Jargon: #1}, coltitle=inkblue, colbacktitle=softgrey, fonttitle=\small
}

% \begin{honest}{Title} ... \end{honest}  - brutally honest finding (rule 18)
\newtcolorbox{honest}[1]{
  enhanced, breakable, colback=white, colframe=warnred,
  boxrule=0.6pt, arc=2pt, left=8pt, right=8pt, top=6pt, bottom=6pt,
  title={\bfseries #1}, coltitle=white, colbacktitle=warnred, fonttitle=\small
}

% \begin{pseudocode}{Caption} ... \end{pseudocode}
% Stands in for algorithm2e, which is NOT installed. Use \Step / \Indent inside.
\newtcolorbox{pseudocode}[1]{
  enhanced, breakable, colback=codebg, colframe=inkblue,
  boxrule=0.4pt, arc=2pt, left=10pt, right=8pt, top=6pt, bottom=6pt,
  fontupper=\ttfamily\small,
  title={\bfseries\rmfamily #1}, coltitle=white, colbacktitle=inkblue, fonttitle=\small
}
\newcommand{\Step}[1]{\par\noindent #1}
\newcommand{\Indent}[1]{\par\noindent\hspace*{2em}#1}

% ===== tikz house styles =================================================
\tikzset{
  box/.style   = {rectangle, rounded corners=2pt, draw=inkblue, fill=softgrey,
                  text=inkblue, align=center, inner sep=6pt, minimum height=9mm,
                  font=\small},
  extbox/.style= {rectangle, rounded corners=2pt, draw=linegrey, fill=white,
                  text=black!70, align=center, inner sep=6pt, minimum height=9mm,
                  font=\small, dashed},
  store/.style = {cylinder, shape border rotate=90, aspect=0.25, draw=accent,
                  fill=white, text=inkblue, align=center, inner sep=5pt,
                  font=\small},
  dec/.style   = {diamond, aspect=1.6, draw=accent, fill=white, text=inkblue,
                  align=center, inner sep=2pt, font=\scriptsize},
  flow/.style  = {-{Stealth[length=2.4mm]}, draw=inkblue, thick},
  dflow/.style = {-{Stealth[length=2.4mm]}, draw=accent, thick, dashed},
  % align=center is load-bearing: without it, a \\ inside an edge label throws
  % the misleading "perhaps a missing \item" error.
  lbl/.style   = {font=\scriptsize\itshape, text=black!65, fill=white,
                  inner sep=2pt, align=center}
}

% forest default for directory trees
\forestset{
  dirtree/.style={
    for tree={
      font=\ttfamily\small, grow'=0, child anchor=west, parent anchor=south,
      anchor=west, calign=first, inner sep=2pt,
      edge path={\noexpand\path [draw, \forestoption{edge}]
        (!u.south west) +(9pt,0) |- (.child anchor) \forestoption{edge label};},
      before typesetting nodes={if n=1{insert before={[,phantom]}}{}},
      fit=band, before computing xy={l=15pt}
    }
  }
}

% ===== misc ==============================================================
\captionsetup{font=small, labelfont={bf,color=inkblue}}
\setlist[itemize]{leftmargin=1.2em, itemsep=2pt, topsep=3pt}
\setlist[enumerate]{leftmargin=1.4em, itemsep=2pt, topsep=3pt}
\newcommand{\file}[1]{\texttt{\small #1}}
\newcommand{\code}[1]{\texttt{\small #1}}
\setcounter{tocdepth}{2}
\setcounter{secnumdepth}{3}


\title{\vspace{-1.5cm}\textbf{Competitive Programming}\\[0.3cm]
\large The Brief}
\author{Explainer for Salman Adnan's portfolio}
\date{}

\begin{document}
\maketitle
\thispagestyle{empty}

\begin{keyidea}{In one paragraph}
361 C++ solutions Salman wrote while training for programming contests (Codeforces, CSES,
AtCoder, CS Academy, and the CS~334 course at Habib University), filed by topic. Three
Python tools make the collection checkable: one compiles every file, one stress-tests the
eight data-structure solutions against brute force, and one tests every solution that was
fixed after a line-by-line review. Thirty files he studied but did not write are kept
apart and never counted. A glossary is at the end.
\end{keyidea}

\tableofcontents

%=====================================================================
\section{The problem the repository solves}
%=====================================================================

Competitive programming is solving short algorithm puzzles against a clock, with an
automated judge running your program on hidden inputs. Salman did it full time for about
eight months, by his own account. The work sat in one untracked folder: duplicates,
empty placeholders, editor backups, half-finished attempts, and reference code from other
people mixed in with his own. The repository turns that folder into something a reader
can browse and a machine can check, and is careful to claim only what it can re-check.

%=====================================================================
\section{What is in it}
%=====================================================================

\begin{figure}[H]
\centering
\begin{forest}
  dirtree
  [competitive-programming/
    [template/\quad\textrm{\small the full contest template, 639 lines}]
    [tools/build\_all.py\quad\textrm{\small compile every file}]
    [stress/stress.py\quad\textrm{\small random inputs against brute force}]
    [solutions/
      [practice/\quad\textrm{\small 297, by topic}]
      [contests/\quad\textrm{\small 41, from 12 live contests}]
      [course-cs334/\quad\textrm{\small 23, from the university course}]
      [from-other-sources/\quad\textrm{\small 30, not Salman's, labelled}]]
  ]
\end{forest}
\caption{The repository at a glance (398 tracked files in all).}
\end{figure}

The practice folders are graphs (65), dynamic programming (36), number theory (48), a
large mixed set (101), and smaller sets on range queries, bitmasks, prefix sums, two
pointers, CS Academy and AtCoder. Folder names say what a training set was meant to
practise, not always what each answer needed: 15 of the graphs files build no graph at
all, and one of them is a sparse table.

%=====================================================================
\section{How a solution is written}
%=====================================================================

A contest program reads from standard input and writes to standard output, nothing else.
Most of Salman's files (267 of 361) paste a shared \emph{template} at the top, then put the
answer in a \code{solve()} function under a \code{SOLVE} banner. Reading any file means
searching for that banner.

\begin{figure}[H]
\centering
\begin{tikzpicture}[node distance=3mm and 6mm]
  \node[box, minimum width=4.2cm] (full) {full library\\105 files};
  \node[box, minimum width=4.2cm, right=of full, fill=accent!20] (trim) {trimmed, about 85 lines\\162 files};
  \node[box, minimum width=4.2cm, right=of trim] (none) {older templates or none\\94 files};
  \node[lbl, below=1mm of full] {April 2025, first};
  \node[lbl, below=1mm of trim] {from mid-April 2025 on};
  \node[lbl, below=1mm of none] {course work, number theory, prefix sums};
  \draw[flow] (full) -- (trim);
\end{tikzpicture}
\caption{Template use across Salman's 361 files. The template shrank over time.}
\end{figure}

The full template (the copy in \file{template/}) is a library of number theory, a sieve,
DSU, a lazy segment tree, a Fenwick tree, graph algorithms, matrix power and fast I/O. File
dates show it came first and was then cut down to a trimmed version keeping only type
aliases, loop macros, gcd, the \code{re}/\code{ps} I/O helpers and debug macros. Local-only
code (timing, reading from a file) is behind a \code{LOCAL == SALMAN} preprocessor guard
that is off on every judge.

%=====================================================================
\section{The core: range queries}
%=====================================================================

The eight CSES Range Queries solutions are the technical heart. Each problem has up to
200{,}000 numbers and 200{,}000 questions, so a loop per question is far too slow; each
structure answers in $O(\log n)$ or $O(1)$.

\begin{table}[H]
\centering\small
\begin{tabular}{@{}p{4.6cm}p{9.6cm}@{}}
\toprule
Problem & Idea \\
\midrule
Range Xor, Forest Queries & prefix tables (1D XOR, 2D sums with inclusion-exclusion) \\
Dynamic Range Sum & Fenwick tree; ``set'' becomes ``add the difference'' \\
Dynamic Range Minimum & segment tree with point updates \\
Hotel Queries & max segment tree that walks down to the leftmost fit \\
Range Update Queries & lazy segment tree, range add \\
Range Updates and Sums & lazy tree with two tags, add and assign, folded so one is pending \\
Polynomial Queries & lazy tag that is a line $c + d\,i$, so progressions combine by addition \\
\bottomrule
\end{tabular}
\end{table}

\begin{figure}[H]
\centering
\begin{tikzpicture}[node distance=8mm and 3mm]
  \node[box] (p) {$[0,3]$: sum $+ 4 \cdot 5$\\\textbf{tag +5}};
  \node[box, below left=of p] (l) {$[0,1]$};
  \node[box, below right=of p] (r) {$[2,3]$};
  \draw[flow] (p) -- (l); \draw[flow] (p) -- (r);
  \node[right=12mm of p, align=left, font=\small] {A covered node updates its own sum and leaves a tag;\\ the tag is pushed to the children only when a later\\ operation needs to look inside. That keeps a range\\ update at $O(\log n)$ instead of $O(n)$.};
\end{tikzpicture}
\caption{Lazy propagation, the idea behind three of the eight solutions.}
\end{figure}

A ninth structure lives outside that folder: a sparse table for static range minimums in
\file{graphs/problem-59}, with $O(1)$ queries from two overlapping power-of-two blocks.

%=====================================================================
\section{Verified, not assumed}
%=====================================================================

\begin{itemize}
\item \textbf{Build check.} \file{tools/build\_all.py} compiles every file in parallel with
  the flags a judge uses, \code{-std=gnu++17} and \code{-O2}. On 26 October 2026, in a clean \code{gcc:13}
  container: 392 of 392 files compile (361 of Salman's, 30 from other sources, the
  template).
\item \textbf{Stress test.} \file{stress/stress.py} generates small random inputs (with
  values up to $10^9$ half the time, to catch overflow), runs the solution and a Python
  brute force, and compares answers token by token, with a fixed seed so failures
  reproduce. All eight range-query solutions matched on 1{,}000 cases each.
\item \textbf{The check can fail.} A deliberately broken minimum query failed on the first
  case.
\item \textbf{Regression tests.} \file{stress/regressions.py} checks each of the 18 solutions
  fixed after review against a brute force written from its problem statement (or a checker,
  or a full game search), built with sanitizers, plus the template's repaired helpers. All
  pass; the old versions of the 12 real bugs fail the same tests.
\item \textbf{One-off sweeps during curation} (not scripts in the repository): every file
  prints an answer, and every file waits for its input instead of crashing at startup (one
  USACO solution reads a named file instead, as USACO requires). Files that failed these
  were removed.
\end{itemize}

%=====================================================================
\section{Code that is not Salman's}
%=====================================================================

A file-by-file authorship pass found 30 files that are not his: 20 reference solutions
headed with Priyansh Agarwal's name (CSES DP, one downloaded folder), and 10 contest
solutions in other people's templates or styles (galen\_colin's and jiangly's templates,
someone else's debug header and input paths, editorial-style code, chatbot-style
comments). Salman edited many of them while learning, so they stay in
\file{from-other-sources/} with a README naming each one, and are excluded from every
count. The first published version counted them; it was corrected within a day.

%=====================================================================
\section{Honest limits}
%=====================================================================

\begin{honest}{What a reader should know}
\begin{itemize}
\item \textbf{What the checks cover.} Correct answers are machine-checked for the eight
  range-query solutions and the 18 fixed ones; for the other 335, the evidence is ``it
  compiles and starts''.
\item \textbf{Bugs that were there.} Reading every file found 12 real bugs (out-of-bounds
  accesses, an operator-precedence slip, an uninitialised counter, a bit test that ignored
  its input, a wrong sort direction, and three solutions whose algorithm was wrong) and 6 latent ones. All are fixed and tested;
  the Deep Dive lists each one.
\item \textbf{Template weight.} The full template's helpers were repaired (a discrete log
  that was off by one, CRT and discrete log that assumed a prime modulus, declarations that
  matched nothing), but it still allocates 41~MB of tables on every run (7~MB for the
  trimmed version).
\item \textbf{Simpler would do.} Range add with point reads uses a lazy segment tree where
  a Fenwick tree on the difference array is a third of the code; the two-tag tree uses
  $-1$ as a sentinel that only works under CSES's limits.
\item \textbf{Mess left in.} Commented-out earlier attempts, unused helpers, a few files
  with numbers but no problem names. None of it prints on a judge.
\end{itemize}
\end{honest}

\section{The story}

\begin{enumerate}
\item About eight months of practice from CS~334 onward (Salman's figure; file dates cover
  April to November 2025, after a bulk copy reset many of them).
\item 25 October 2026: curated from 480 candidate files into a repository; duplicates,
  placeholders, backups and non-compiling attempts left out, everything compiled, the
  range queries stress-tested, an assistant-written ICPC 2026 folder excluded.
\item 26 October 2026: a second, file-by-file pass corrected the template history, moved
  30 files by other people into their own folder, and dropped seven non-solutions. The
  count settled at 361.
\item Later the same day: every catalogue line was checked against its file, and the bugs
  that review found were fixed, each with a brute-force test.
\end{enumerate}

\textbf{If starting over:} write the brute force first, keep helpers in one header,
prefer the simplest structure that fits, and use version control from day one.

%=====================================================================
\section{Glossary}
%=====================================================================

\begin{description}[leftmargin=0pt, itemsep=3pt]
\item[Brute force] the slow, obviously correct way to compute an answer; the reference a
  stress test compares against.
\item[CSES] a fixed set of 300 training problems by topic.
\item[Fenwick tree] an array of partial sums indexed by binary structure; point add and
  prefix sum in $O(\log n)$.
\item[Lazy propagation] deferring a range update by leaving a tag on a covering node.
\item[$O(\log n)$] grows very slowly with input size: about 18 steps for 200{,}000.
\item[Segment tree] a binary tree of range summaries; any range is a handful of nodes.
\item[Sparse table] precomputed minimums of power-of-two blocks; $O(1)$ static queries.
\item[Stress test] many random small inputs, fast answer against brute force.
\item[Template] shared helper code pasted at the top of every solution.
\item[Upsolving] solving a problem after the contest, often from the editorial or others'
  code.
\end{description}

\end{document}
